\section{Introduction}

A combinatorial circuit is any circuit that maps inputs to outputs without retaining internal state.
These circuits have many important application, such as in complexity theory~\cite{Vollmer1982}, hardware design~\cite{HarrisWeste2010}, and quantum program compilation~\cite{CrossJavadiAbhari2022}.
This wide range of applications has motivated the formal mathematical study of combinatorial circuits.
One early approach to this problem was to view combinatorial circuits as string diagrams in a monoidal category~\cite{Hotz1965}.
This approach has seen wide acceptance in the field of quantum program design and analysis~\cite{HeunenVicary2019}.
In this model, each gate is a morphism $G: n \to m$ with $n$ input wires and $m$ output wires.
If the number of outputs from a gate $G$ matches the number of inputs to a gate $H$, then $H \circ G$ denotes the circuit obtained by connecting the outputs of $G$ to the inputs of $H$.
Given any two gates $G$ and $H$, $G \otimes H$ denotes the circuit obtained by placing $G$ and $H$ side-by-side without any of their inputs connecting.
In other words, $( \circ )$ denotes sequential composition and $( \otimes )$ denotes parallel composition.
Given a set $\Sigma$ of gates, $F( \Sigma )$ denotes the family of circuits obtained from the gates in $\Sigma$.
These circuits form a \emph{free monoidal category}.
Their semantic interpretation is then given by a \emph{monoidal functor} $\sem{-}: F( \Sigma ) \to \mathcal{D}$ where $\mathcal{D}$ denotes the semantic domain.
It is easy to show that $\sem{-}$ is fully determined by where it sends the gates in $\Sigma$, as is expected of combinatorial circuits.
This perspective on combinatorial circuits has had many successes, such as providing decision procedures for the equality of certain classes of classical circuits~\cite{Lafont2003}, and optimization procedures for certain classes of quantum circuits~(see~\cite{Backens2021} for an introduction).

In most quantum circuit semantics, $\mathcal{D}$ is taken to be the category of linear transformations between some class of spaces.
However, with the advent of quantum machine learning, this approach is no longer sufficient.
As in classical machine learning, a quantum machine learning model is a tensor network (i.e., a circuit) with one or more free parameters from some continuous space~(see~\cite{HugginsPatil2019}).
To see why this is challenging, consider a simple quantum machine learning model with two parameters $\{ p_1, p_2 \}$ which can be used to instantiate Pauli-$Z$ rotations denoted by the gate $R_Z( - ): 1 \to 1$.
For example, $R_Z( p_1 )$ and $R_Z( p_1 + p_2 )$ are both instantiations of the gate $R_Z( - )$.
Following the standard semantics, $\sem{R_Z( p_1 )}$ and $\sem{R_Z( p_1 + p_2)}$ should correspond to the following continuous maps with domain $\mathbb{R}^2$.
\begin{align*}
    \sem{R_Z( p_1 )}( x, y ) &=
    \begin{bmatrix}
        e^{-ix/2} & 0 \\
        0 & e^{ix/2}
    \end{bmatrix}
    &
    \sem{R_Z( p_1 + p_2 )}( x, y ) &=
    \begin{bmatrix}
        e^{-i(x+y)/2} & 0 \\
        0 & e^{i(x+y)/2}
    \end{bmatrix}
\end{align*}
One would then expect that $R_Z( p_1 )$ and $R_Z( p_1 + p_2 )$ compose as follows.
\begin{equation*}
    \sem{R_Z( p_1 ) \star R_Z( p_1 + p_2 )}( x, y )
    =
    \begin{bmatrix}
        e^{-ix/2} & 0 \\
        0 & e^{ix/2}
    \end{bmatrix}
    \begin{bmatrix}
        e^{-i(x+y)/2} & 0 \\
        0 & e^{i(x+y)/2}
    \end{bmatrix}
\end{equation*}
It should be noted that this new composition is neither the composition of continuous functions, nor the composition of linear functions.
Instead, this composition is best understood as: (1) copying the parameters; (2) providing a copy of the parameters to each function; (3) composing the resulting linear maps.
This composition has appeared ad-hoc in many prior works~\cite{JeandelPerdrix2018,MillerBakewell2020,XuLi2022,PehamBurgholzer2023,HongHuang2024}.
In an extended abstract~\cite{Wesley2025}, we noted that vector spaces are enriched over topological spaces, and that this enrichment explains why the construction is well-defined.
However, the underlying properties of enriched categories which make this construction possible have yet to be studied.

One alternative approach to parameterized quantum circuits is that of linear dependent type theory.
In~\cite{FuKishida2022}, it was shown that set-theoretic parameters could be introduced to a quantum programming language without violating the linearity of the language.
Moreover, the parameters could be used to parameterize not only the family of morphisms, but also their domains and codomains.
However, these parameters were introduced freely, meaning that the set of parameters would have no internal structure.
This means that the semantic model could not enforce that every quantum circuit is a continuous function, as required by quantum machine learning.
Moreover, when the parameter spaces are taken to be topological spaces, then the topology of the spaces can be used to encode additional functional requirements through the type system.
For example, if a gate $G( p )$ is parameterized by the space $\mathbb{R} / (2\pi \mathbb{Z} )$, then $\sem{G( p )}( x )$ must have period of at most $2\pi$, which cannot be enforced through the types in a purely set-theoretic model of parameters.

More fundamentally, the model of linear dependent type theory given in~\cite{FuKishida2022} does not account for the new sequential composition found in the motivating example.
In principle, the parameters could be copied and passed along each time two parameterized gates are composed.
However, given a fixed set of named parameters, there should never exist a case in which the composition is performed without first copying the parameters.
For this reason, the present paper looks for a new notion of parameterization which encodes the natural choices for sequential and parallel composition.
Since the parameterized families of maps go from parameter objects in some category $\mathcal{V}$ to hom-sets in some category $\mathcal{C}$, then it is clear that $\mathcal{C}$ must be enriched over $\mathcal{V}$.
Given this enrichment, it can then be asked what is necessary for the new sequential and parallel composition to be well-defined.
These questions turn out to be interesting not only from the perspective of parameterized quantum circuits, but also from the perspective of enriched category theory.

This new notion of parameterization also provides new insights into traditionally well-understood constructions in quantum computing.
For example, controlled operations can be understood as an instance of this new notion of parameterization.
Intuitively, the control qubits can be thought of as linear parameters, which are used to multiplex between the operations performed on the target qubits.
Since the control qubits are not modified during this procedure, then in some sense they are never consumed, and will persist after the controlled operation is finished.
Alternatively, we can think of the control qubit being copied in some linear way, with only the copy being consumed during the operation.
This new point-of-view is not only insightful, but allows for many well-known circuit relations to be re-obtained in a more structured manner.
This construction is in line with the comonoidal view of quantum control found in~\cite{HeunenVicary2019}.

The rest of this paper is structured as follows:
\cref{Sect:Background} briefly reviews the background required to read this paper;
\cref{Sect:Motivation} provides motivation for the paper through the study of quantum circuits with parameterized rotations (e.g., ansatz circuits);
\cref{Sect:Construction} outlines a generalized construction for parameterized categories, which is shown to have nice categorical properties and immediate applications to quantum information;
\cref{Sect:CartMon} studies the special case of parameterization over Cartesian monoidal categories, in which all objects in $\mathcal{V}$ can be seen as classical parameters (e.g., they can be copied and discarded);
\cref{Sect:Closed} studies the special case of self-enrichment as a generalization of control functors.
